\subsection{Proper fusion subrings and Galois reduction}

\begin{proposition}\label{prop:rank4-galois-exception}
Let $R$ be a self-dual rank-four fusion ring. Unless
$R\otimes_{\mathbb Z}\mathbb Q\cong K\times K$ for a real quadratic
field $K$, every pivotal categorification of $R$ has a pseudo-unitary
Galois conjugate, or is already pseudo-unitary.
\end{proposition}
\begin{proof}
The rational algebra is a product of totally real number fields whose
degrees sum to four. Let $K_0$ be the field corresponding to the
Frobenius--Perron character and put $D=\operatorname{FPdim}(R)$.
If $K_0=\mathbb Q$, the category is integral and pseudo-unitary. If
$D\in\mathbb Q$, it is weakly integral and again pseudo-unitary.
Assume therefore that both are nonrational.

A pivotal categorification is spherical. Its dimension field must contain
$D$. If $[K_0:\mathbb Q]=4$, every character belongs to the same Galois
orbit. If its degree is three, the only other field factor is rational,
so cannot supply the spherical dimension character. If its degree is
two, then $D$ generates $K_0$. Every other field factor is rational,
or there is one further quadratic factor $K_1$. Rational factors cannot
contain $D$, and a distinct quadratic field $K_1$ cannot contain it either.
Except in the indicated repeated-field case, the dimension character
therefore belongs to the Frobenius--Perron orbit. Galois conjugation of
the spherical category gives the claimed pseudo-unitary realization.
\end{proof}

The repeated-field exception cannot be removed for individual realizations:
the product of a Fibonacci category and its Yang--Lee Galois conjugate has
rank four and categorical global dimension $5$, whereas the Frobenius--Perron
dimension of its fusion ring is $5\phi^2$, with $\phi=(1+\sqrt5)/2$.
Its global dimension stays $5$ under Galois conjugation. Its fusion ring
nevertheless has a unitary realization, namely the product of two unitary
Fibonacci categories.

\begin{theorem}\label{thm:proper-subring-unbounded}
A rank-four fusion ring with a proper nontrivial fusion subring admits a
pivotal categorification if and only if it is one of the thirteen rings
in the realization table other than $F_{1,6}$ and $F_{2,15}$.
All thirteen admit unitary categorifications. There is no multiplicity bound.
\end{theorem}
\begin{proof}
If the ring has a nontrivial invertible basis element,
Corollary~\ref{cor:nontrivial-invertible-unbounded} gives precisely eleven
possibilities. Every such ring has a proper pointed subring (or, in the
pointed rank-four cases, a subgroup of order two).

Suppose the pointed part is trivial. The unrestricted non-self-dual
classification implies that a pivotal-categorifiable ring of this form
is self-dual. A proper nontrivial fusion subcategory has rank two or
three and inherits the pivotal structure. In rank two, the
moment obstruction~\cite[Proposition~\ref{B-prop:ranktwo}]{Finite} leaves the Fibonacci fusion ring.
In rank three, the pivotal classification~\cite{OstrikRank3} leaves
only the ring of $\operatorname{PSU}(2)_5$, because every other positive
rank-three ring has a nontrivial invertible object. But the single
remaining simple $W$ outside such a rank-three subcategory must satisfy
$VW=n_VW$ for every simple $V$ in it: reciprocity excludes summands in
the subcategory from $VW$. Thus $\operatorname{FPdim}(V)=n_V\in\mathbb Z$,
contradicting the nonintegral dimensions of $\operatorname{PSU}(2)_5$.
A Fibonacci subring is therefore present.
Theorem~\ref{thm:unbounded-fibonacci} leaves exactly
$\mathrm{Fib}\boxtimes\mathrm{Fib}$ and $\operatorname{Ad}(E_8)$
when the pointed part is trivial. Together with the eleven preceding
rings these are the stated thirteen. The displayed unitary models prove
sufficiency.
\end{proof}

